\documentclass[11pt]{article}

\usepackage[margin=0.9in]{geometry}
\usepackage[T1]{fontenc}
\usepackage{lmodern}
\usepackage{microtype}
\usepackage{booktabs}
\usepackage{enumitem}
\usepackage{amsmath}
\usepackage[hidelinks]{hyperref}

\setlist{itemsep=2pt, topsep=4pt}
\newcommand{\Walker}{Hangarback Walker}

\title{Hangarback Walker in Three-Card Blind:\\
  why the Elixir player wins on the play and draws on the draw}
\author{}
\date{November 2026}

\begin{document}
\maketitle

\begin{abstract}
\noindent
Player A has two \Walker{}s and City of Traitors. Player B has one
\Walker, Elixir of Immortality and City of Traitors. This matchup was
judged too complex to resolve, and \Walker{} was banned from Three-Card
Blind. We settle it: when B plays first, B wins; when A plays first,
neither player can force a win, so the game is a draw. Both results are
theorems in the Rocq proof assistant. This note explains in terms of the
cards why they hold, and sketches how they were proved.
\end{abstract}

\section{The answer}

Three-Card Blind games are played with perfect information and optimal
play. Both players start on 20 life with their three cards in hand. Drawing
from an empty library does nothing. Random effects go the opponent's way,
so when Elixir shuffles B's library, A chooses the order.

The clearest way to see the results is to ask how much life B would need.
For each opening there is a level of B's starting life at or below which
A can force a win, and a level at or above which B can force a win:

\begin{center}
\begin{tabular}{lcc}
\toprule
Who plays first & A wins if B starts on & B wins if B starts on \\
\midrule
B & 10 or less$^{\dagger}$ & 17 or more$^{*}$ \\
A & 15 or less$^{\dagger}$ & 24 or more$^{\dagger}$ \\
\bottomrule
\end{tabular}
\end{center}

\noindent
B actually starts on 20. On the play, 20 is in B's winning range, and
Rocq proves that B wins. On the draw, 20 falls in the gap between the two
levels, and Rocq proves both halves of the draw: A can't win from 20, and
B can't win from 20 either. Being on the play moves both levels in B's
favour; \S\ref{sec:tempo} shows where that comes from.

{\small
\noindent
$^{*}$ B's proof strategy needs only 17 of its 20 life.
$^{\dagger}$ Found by the solver with the same sound methods as the
proofs, but not checked in Rocq. Each is a bound: A wins at least up to
the number shown, and B wins from at least the number shown. That these
are the \emph{exact} levels is the solver's best estimate.\par}

\section{Five facts that drive every game}
\label{sec:facts}

\begin{enumerate}[label=\textbf{F\arabic*.}, leftmargin=*]
\item \textbf{Two mana a round.} City untaps only in its owner's untap
  step, and mana empties between steps. A player spends their 2 mana
  either on their own turn (to cast) or by leaving City untapped and
  using it on the opponent's turn (to pump or activate Elixir).
\item \textbf{A Walker can attack or grow in a round, but not both.} A
  pump costs \{1\},\{T\}, so a tapped Walker can't be pumped. A Walker
  that attacks stays tapped through the opponent's turn. A Walker pumped
  on its owner's turn is also still tapped on the opponent's turn, so it
  can't block. Players therefore pump on the opponent's turn, after that
  turn's combat, and keep their Walkers untapped as blockers until then.
  B also pumps in its own second main phase when 1 mana is left over
  after casting Elixir. A new Walker is summoning sick until its owner's
  next turn, so it can't be pumped before then.
\item \textbf{Killing a Walker arms its owner.} A dying Walker leaves one
  flying 1/1 Thopter per counter. Walkers can't block fliers; only
  Thopters can. So the side that wins a Walker fight hands the loser an
  evasive army of the same size.
\item \textbf{Elixir is a 3-mana loop.} Casting costs 1 and activating
  costs 2. Activating gains 5 life and shuffles Elixir and B's graveyard,
  including a dead Walker, back into B's library. B then draws them again
  on later turns. (In every line shown here A puts Elixir on top, so a
  dead Walker comes back one draw later. The proofs cover both orders.)
  At 2 mana a round, B can gain 5 life about every two rounds.
\item \textbf{A's losses are permanent.} A has no recursion and no life
  gain. A dead A Walker stays dead.
\end{enumerate}

\noindent
So the decks pull in opposite directions. \emph{A has growth}: two
Walkers, 2 counters a round while they stay home. \emph{B has life and
recursion}: 5 life every two rounds, and a Walker that always comes back.
Every combat decision weighs one against the other.

\section{B on the play: B wins}

\paragraph{The opening.} The proof's strategy, against an A that keeps
its Walkers home:

\begin{center}
\begin{tabular}{clll}
\toprule
Round & B's turn & A's turn & Life (A, B) \\
\midrule
1 & Walker, $X=1$ & Walker, $X=1$ & 20, 20 \\
2 & pump to 2/2, cast Elixir & second Walker, $X=1$ & 20, 20 \\
3 & 2/2 attacks, A doesn't block & --- & 18, 20 \\
4 & 2/2 attacks, A doesn't block; & B activates Elixir & 16, 25 \\
  & \quad A pumps both Walkers to 2/2 & \quad in A's end step & \\
\bottomrule
\end{tabular}
\end{center}

\noindent
Going first, B has a 2/2 attacking in round 3, while A has just spent
both of its first two rounds of mana casting two 1/1s. Why doesn't A
block? A double block kills the 2/2, but B can assign its 2 damage to kill
both 1/1s. A then has two Thopters and no Walkers, for good. B has two
Thopters, and its Walker comes back with the next Elixir.

\paragraph{A's dilemma.} From here, B's 2/2 or 3/3 attacks every turn
it can, even into bigger untapped A Walkers. Both of A's answers lose
ground:
\begin{itemize}
\item \emph{Block.} B's Walker dies into Thopters that A's Walkers can't
  block (F3), and Elixir brings the Walker back (F4). A swaps damage on
  the ground for damage in the air, and B has lost nothing for good.
\item \emph{Don't block.} A takes 2 or 3 a turn.
\end{itemize}
A can't race either. Every A Walker that attacks gives up a counter and
a blocker for the round (F2), while Elixir hands B back 5 life every two
rounds and nothing hands A back anything (F5). In the proof's main line
A first blocks on B's ninth turn, with a 6/6. B's 3/3 becomes three
Thopters, which finish A in the air.

\paragraph{The game ends.} B's strategy always makes progress. From any
position it reaches, B deals A another point of damage within a bounded
number of turns (17 from the opening), and A has only 20 life to lose.
The longest game this allows is under 800 turns; sample lines end in 19
to 55.

\paragraph{The Elixir is the whole matchup.} If B is not allowed to use
Elixir, the solver finds no B win from any starting life up to 100. It
also finds no way for B to avoid losing when A goes first. (These are
solver checks with size limits, not theorems.)

\section{A on the play: a draw}

A draw means neither side can force the other to 0. With best play the
game goes on forever: B keeps gaining life and A keeps growing, and
neither can break through.

\subsection{What changes on the draw}
\label{sec:tempo}
For two rounds the openings look alike. Both players cast a Walker; then
A casts its second Walker while B pumps to 2/2 and casts Elixir. What
changes is where B's first attack falls in the turn order.
\begin{itemize}
\item \emph{On the play,} B's first attack (round 3) comes straight after
  A's turn 2, in which A spent City's mana on its second Walker. A is
  tapped out and can't grow anything. The 2/2 hits again in round 4,
  and A's Walkers reach 2/2 only after that, with A on 16.
\item \emph{On the draw,} A's turn 3 comes before B's first attack. A
  casts nothing and keeps City untapped. Once the 2/2 has hit, A pumps
  both Walkers to 2/2, with A on 18. In the solver's line, B's 2/2 stays
  home on its next turn, facing two untapped 2/2s.
\end{itemize}
The solver puts a number on that one round of mana. Take B's turn with
its first attack, B's 2/2 and Elixir against two 1/1s. B needs 16 life to
win when A is tapped out with its second Walker just cast, and 24 when A has City's mana and both Walkers
ready$^{\dagger}$. The same solver table gives 17 for the B-first opening
and 24 for the A-first one. So A's mana during B's first attack accounts
for the whole difference between the openings.

\subsection{A can't win}
B is safe from 20 life. B plays much as on the play, but keeps home
just enough Thopters to chump-block, so that A's all-out attack next
turn is not lethal, counting Elixir's 5 life. When an attack is lethal,
B chumps A's biggest Walkers first and activates Elixir after blocks.

The proof's main line shows A racing. From turn 14, A
attacks every other turn with a 6/6 and a 4/4, 10 damage a swing, while
B's 4/4 hits back. B takes $4+10+10+10 = 34$ damage, gains 15 from three
Elixirs, and ends on $20+15-34 = 1$ life. Then B's next swing kills A.
Without Elixir, B would have died to A's second big swing.

\subsection{B can't win}
A is safe as long as B has at most 20 life. Roughly, A plays like this:
\begin{itemize}
\item On its own turn A only casts Walkers. City's mana is for B's turn,
  where A pumps in the end step.
\item A's Walkers attack only when each is bigger than B's Walker (or B has none). A
  second case: if B's Walker is tapped, A sends its bigger Walker alone
  (once it has 3 counters) and keeps the smaller home to grow.
\item A blocks only a large B Walker that is bigger than all of its own,
  and uses its smallest Walker to do it.
\end{itemize}

\paragraph{Why A keeps up with Elixir.} B can gain 5 life every two
rounds forever, so no fixed amount of damage keeps A safe. A's threat has
to grow too, and it can. Elixir and growth compete for B's mana. Each
loop costs B 3 mana for 5 life, plus 2 more to recast a dead Walker. A
spends all of its mana on growth, and once both its Walkers are 3/3 it
puts every counter on the bigger one. The proof's main line shows the
effect. Here B picks, every turn, the move that leaves A the least room.
\begin{itemize}
\item \emph{Midgame.} A's 3/3 attacks whenever B's Walker is tapped,
  and A's other Walker grows one counter a round. By turn 16 it is a 7/7,
  and B's Walker has stopped at 4/4.
\item \emph{The trade.} B's 4/4 attacks; A blocks with its 3/3 and
  pumps it, so both die and each side gets four Thopters. The Thopters
  trade off on A's next attack.
\item \emph{Endgame.} A's lone 7/7 is left against B's Elixir loop.
  B's Walker comes back each time as a 1/1, good only for chump-blocking,
  and a loop that recasts it takes B about two and a half rounds of mana.
  The 7/7 attacks every turn. B chumps some attacks and gains 5 a loop,
  but enough 7-point hits get through. B's life runs 15, 8, 1, 6, 11, 4,
  9, 2, 7, and then 0. B's last Thopter pings A for 1 a turn, from 11
  down to 3, which is too slow.
\end{itemize}
The proof keeps score with the most life B can have while A stays safe.
Along this line that number equals B's actual life at every turn: each
Elixir raises both by 5, and each hit lowers both. B never gets ahead.
B can avoid this line by playing for safety instead, so the result is a
draw, not a win for A.

\paragraph{A must hold Thopters back.} Once Walkers have traded into
Thopters, A has to split its Thopters, some attacking and some at home
to block B's Walker. A proof that lets A attack only with all of its
Thopters or none can't reach 20 (even an optimistic version of that check
stops at 19). Allowing A to keep up to two Thopters home reaches exactly
20.

\subsection{Why exactly 20?}
\label{sec:open}
Both halves of the draw have no margin. B is safe from exactly 20 life,
and A is safe up to exactly 20. That comes from the proofs, not the game.
Each proof gives its own side fewer options than the real game has
(\S\ref{sec:proof}), and those limits were loosened only until each
proof reached 20. With looser limits the solver puts the levels at 15
for A and 24 for B, so 20 sits well inside the drawn range.

\section{How it was proved}
\label{sec:proof}

\paragraph{The model.} The rules are the Comprehensive Rules in force
when \Walker{} was banned, with damage assignment order, plus the
Three-Card Blind format rules. The model is written in Rocq. It leaves
out only what can't matter with these six cards, and each omission is
justified in writing. For example, there is no stack: nothing here
targets or counters anything, and the effects give the same result in any
order, so responding never changes the outcome. Rocq also proves that
the model never gets stuck, never creates or loses a card, and has no
move that does nothing.

\paragraph{What is claimed.} ``B wins'' means B can force A to 0 life in
finitely many moves, whatever A does. ``A is safe'' means A has a set of
positions that contains the start and no lost position, and that B can
never force A out of. A draw means both players are safe.

\paragraph{Making an infinite game finite.} Walker counters, Thopters and
B's life are all unbounded. A \emph{certificate} is a finite table of
positions, one for each turn start, that a checker can verify one turn at
a time. Five ideas keep it finite:
\begin{itemize}
\item \emph{Only turn starts are stored.} The play within each turn is
  searched in full.
\item \emph{B's life is not part of the position.} More life never hurts
  B, so each entry stores only the least life B needs, or the most life A
  can allow.
\item \emph{The side being proved may be restricted.} Giving up some of
  your own moves can't create a win you don't have. For example, B's
  winning strategy never grows its Walker past 4 counters.
\item \emph{The opponent gets every legal move,} but its large numbers
  are cut off in a way that can only hurt the side being proved. An
  opponent Walker past a size limit counts as being of unknown size, and
  any event where its exact size would matter counts as a loss.
\item \emph{A's growing threat.} For A's safety, a size limit on A's own
  Walkers can't work, because B's life grows without bound. Instead, one
  A Walker past the limit stands for every larger size at once, $K+n$
  counters for all $n \ge 0$. Each entry then stores the most life B may
  have as a linear function of $n$.
\end{itemize}

\paragraph{Checking.} An untrusted solver computes the tables. Rocq
proves, once and for all, that any table whose entries all pass the
checker proves the claim. Then it checks every entry: about 0.57 million
positions for the B-first win, and 0.13 million and 0.58 million for the
two halves of the A-first draw, in about 40 CPU-minutes. A wrong table
can only make the check fail. Only Rocq itself has to be trusted.

\section{What is not settled}
\begin{itemize}
\item \textbf{Short rule strategies are summaries.} Rules like A's in
  \S4.2, written for both players, agree with the proofs' strategies at about 98.5--99.4\% of decisions. B's
  most common exception is attacking with the Walker while Elixir is in
  hand. With no exceptions, the rule strategies fail the solver's checks.
  Patching B's exceptions one at a time made the total worse. Many
  exceptions have no game reason we could find, and they may come from
  the certificates' limits rather than from the game.
\item \textbf{The exact levels.} The $\dagger$ numbers in \S1 have not
  been checked in Rocq. Whether B could also win on the play from 11--16
  life is open.
\item \textbf{The value of tempo.} The 8 life that A's held mana is worth
  (\S\ref{sec:tempo}) comes from the solver, not from the proofs.
\end{itemize}

\noindent
The model, proofs and solver are at\\
\url{https://github.com/Skyb0rg007/three-card-blind-proofs}.

\end{document}
